% !TeX encoding = UTF-8
\chapter*{Abstract}

In the fourth lab of the Security Insider Lab, we build a complete real-time DDoS defense cycle with eBPF/XDP in which the classifier and the filter run inside the Linux kernel. One XDP program captures packets at the network driver hook, computes ten forward-direction flow features per flow in a kernel map, classifies the flow with a Decision Tree (depth 6) and an XGBoost model (20 trees, depth 4) stored in BPF maps, and drops all packets of a source that is classified as malicious. User space only trains and exports the models, loads and configures the program, and displays results. On the 33,463 test flows of the CIC-IDS2017 Friday DDoS file, the integer models reach an F1 score of 0.9991 (Decision Tree, 4 false alarms, 32 missed attacks) and 0.9989 (XGBoost, 11 false alarms, 31 missed attacks). Their verdicts agree with the floating-point models on all test flows, and the kernel code reproduces the Python integer evaluator exactly on 1,024 test vectors. In a Docker topology with an inline IDS gateway in front of Nginx and Apache, normal client traffic produced no malicious verdict and no drop, and a stub model that always answers malicious showed the filter dropping 55 packets of the normal client. Without attack tools, 60 dataset flows replayed through the attached kernel program were handled correctly: all 30 DDoS flows were classified malicious and blocked, and none of the 30 benign flows. We also configured xdp-filter and xdpdump (Task 1), built a web interface (Task 3) and containerized Nginx, Apache and a pinging container (Task 4). The models are essentially a payload-size rule, which we discuss together with the other limits of the in-kernel design.
